\section{Tema 4}
\section*{A) EL SECTOR ENERGÉTICO}

\begin{enumerate}[label=\textbf{\arabic*.}]

    \item \textbf{El sector energético español se caracteriza por un elevado grado de internacionalización de sus empresas líderes, especialmente en mercados latinoamericanos.} \\
    \textbf{Verdadera.} \textit{Justificación: Como se señala en la página 3, uno de los rasgos específicos del sector es que existe un elevado grado de internacionalización de las principales empresas españolas en países como Perú y Venezuela.}

    \item \textbf{La industria es el sector económico que menor cantidad de energía consume, actuando principalmente como un proveedor de inputs para los servicios.} \\
    \textbf{Falsa.} \textit{Justificación: En la página 4 se indica exactamente lo contrario: la energía es un input fundamental para el resto de sectores, ``en especial, para la industria que es la que mayor consume''.}

    \item \textbf{El proceso de privatización de grandes empresas energéticas en las últimas décadas ha eliminado las estructuras de oligopolio en el mercado español de hidrocarburos.} \\
    \textbf{Falsa.} \textit{Justificación: Según la página 11, a pesar de la privatización, se ha generado un elevado grado de concentración. En hidrocarburos, Repsol, Cepsa y BP conforman un oligopolio que controla más del 70\% del mercado.}

    \item \textbf{El incremento del consumo energético en las actividades productivas españolas se explica, en parte, por el proceso de sustitución de mano de obra por energía.} \\
    \textbf{Verdadera.} \textit{Justificación: La página 5 explica que el gran crecimiento del consumo en las actividades productivas se da ``como resultado de la sustitución de mano de obra por energía''.}

    \item \textbf{España ha alcanzado históricamente el equilibrio en su balance energético gracias a un elevado grado de autoabastecimiento basado en recursos propios.} \\
    \textbf{Falsa.} \textit{Justificación: La página 8 advierte que el balance energético español registra históricamente un desajuste, ya que la producción nacional es insuficiente y el grado de autoabastecimiento es de ``sólo del 30\% de las necesidades de energía primaria''.}

    \item \textbf{En la fase de transporte de electricidad en alta tensión, el mercado opera en régimen de libre competencia entre Red Eléctrica Española e Iberdrola.} \\
    \textbf{Falsa.} \textit{Justificación: Según la página 11, en la fase de transporte en alta tensión no hay libre competencia, sino que ``el monopolio lo ostenta Red Eléctrica Española''.}

    \item \textbf{El sector energético se considera estratégico no solo por su valor propio, sino por su capacidad de generar efectos de arrastre que mejoran la eficiencia del resto de actividades productivas al ser un input fundamental.} \\
    \textbf{Verdadera.} \textit{Justificación: La página 4 confirma que el sector energético tiene ``efectos de arrastre sobre otras ramas, de forma que si el sector gana en eficiencia, se benefician el resto de actividades productivas''.}

    \item \textbf{Gracias al reciente e intenso crecimiento de las energías renovables, la producción nacional de energía en España es ya suficiente para satisfacer la totalidad de la demanda interna de las actividades productivas.} \\
    \textbf{Falsa.} \textit{Justificación: Aunque las renovables han crecido, la página 10 establece como característica principal del sector que existe una ``producción nacional insuficiente para satisfacer la demanda'' y una ``fuerte dependencia del exterior''.}

    \item \textbf{En el sistema eléctrico español, las empresas dominantes mantienen una integración vertical total que incluye el control absoluto de la fase de transporte en alta tensión.} \\
    \textbf{Falsa.} \textit{Justificación: La página 11 indica que ``la integración vertical es parcial'' en generación, distribución y comercialización, pero el transporte de alta tensión queda fuera de su control, al ser monopolio de Red Eléctrica Española.}

    \item \textbf{La transformación de la estructura de la demanda energética en España durante la década de los noventa se caracterizó principalmente por el proceso de sustitución del petróleo por el gas natural.} \\
    \textbf{Verdadera.} \textit{Justificación: En la página 5 se detalla la cronología de la estructura de la demanda, indicando explícitamente que en los años 90 se produjo la ``sustitución de petróleo por gas natural''.}

    \item \textbf{Tras los procesos de privatización de grandes compañías como Repsol o Endesa, el sector público español ha retirado por completo su participación en todas las ramas de actividad del sector energético.} \\
    \textbf{Falsa.} \textit{Justificación: La página 10 aclara que el proceso de privatización ha reducido de forma notable la participación pública, ``exceptuando el sector del carbón'', por lo que no se ha retirado por completo.}

\end{enumerate}

\section*{B) EL SECTOR INDUSTRIAL MANUFACTURERO}

\subsection*{1) DELIMITACIÓN Y CLASIFICACIÓN}

\begin{enumerate}[label=\textbf{\arabic*.}]
    \setcounter{enumi}{11}

    \item \textbf{La industria manufacturera se distingue por tener una productividad del trabajo superior a la de la agricultura y los servicios, incorporando más rápidamente el progreso técnico.} \\
    \textbf{Verdadera.} \textit{Justificación: La página 14 señala entre los rasgos característicos que la industria tiene mayor productividad laboral al ser ``más susceptible de incorporar el progreso técnico, ahorrador de mano de obra''.}

    \item \textbf{Siguiendo los criterios del Sistema Europeo de Cuentas (SEC), el concepto de industria incluye tanto las actividades extractivas como todas las manufactureras.} \\
    \textbf{Falsa.} \textit{Justificación: La página 15 indica expresamente que ``El SEC considera industria sólo a las actividades manufactureras'', excluyendo las extractivas.}

    \item \textbf{Las ramas industriales de ``demanda fuerte'' son aquellas vinculadas a sectores tradicionales como el textil y la siderurgia debido a su estabilidad histórica.} \\
    \textbf{Falsa.} \textit{Justificación: En la página 15 se clasifica al textil y la siderurgia como ramas de ``Demanda débil''. Las de demanda fuerte son las de reciente desarrollo y mayor innovación, como la aeroespacial y electrónica.}

    \item \textbf{Las empresas pertenecientes a sectores industriales avanzados operan generalmente en mercados internacionales muy competitivos y presentan una oferta más concentrada.} \\
    \textbf{Verdadera.} \textit{Justificación: La página 16 describe a las industrias avanzadas afirmando que tienen una ``Oferta más concentrada en las empresas líderes'' y un ``Mayor nivel de apertura exterior, actuando en un mercado internacional más competitivo''.}

    \item \textbf{Las actividades industriales tradicionales se caracterizan por emplear mayoritariamente mano de obra con alta cualificación universitaria en comparación con las avanzadas.} \\
    \textbf{Falsa.} \textit{Justificación: Según la página 16, son las ramas avanzadas (y no las tradicionales) las que se caracterizan porque ``Utilizan mano de obra más cualificada''.}

\end{enumerate}

\subsection*{2) EVOLUCIÓN DEL SECTOR}

\begin{enumerate}[label=\textbf{\arabic*.}]
    \setcounter{enumi}{16}

    \item \textbf{La pérdida de peso relativo de la manufactura en el VAB y el empleo es una tendencia que España comparte con las principales economías europeas.} \\
    \textbf{Verdadera.} \textit{Justificación: La página 20 confirma la pérdida de peso de la industria, añadiendo que ``Esta tendencia es compartida con las principales economías europeas''.}

    \item \textbf{El proceso de desindustrialización en España ha sido más severo si se mide en términos de VAB real que si se analiza a través de la evolución del empleo.} \\
    \textbf{Falsa.} \textit{Justificación: La página 20 afirma lo contrario: ``La contracción en el peso de las manufacturas en el VAB real es menor que respecto al VAB a precios corrientes y al empleo''.}

    \item \textbf{La progresiva externalización de servicios digitales y TIC por parte de la industria ha provocado un aumento estadístico del peso del sector industrial en el conjunto de la economía.} \\
    \textbf{Falsa.} \textit{Justificación: En la página 20 se expone que la progresiva externalización de servicios es una de las razones que explican la reducción o pérdida de peso relativo del sector industrial frente a los servicios, no un aumento estadístico.}

    \item \textbf{La producción industrial española muestra una intensidad de fluctuación cíclica superior a la del conjunto de la economía nacional.} \\
    \textbf{Verdadera.} \textit{Justificación: La página 21 afirma que existe una ``Mayor intensidad de las fluctuaciones cíclicas en la industria española'', retrocediendo más en depresiones y avanzando más en recuperaciones.}

    \item \textbf{Tras la crisis sanitaria de la Covid-19, la industria española ha logrado recuperar y superar su cuota de participación en la producción industrial total de la Unión Europea.} \\
    \textbf{Falsa.} \textit{Justificación: La página 23 señala explícitamente que ``Tras la pandemia de la Covid-19, se volvió a reducir la participación de la industria española en la producción industrial europea''.}

    \item \textbf{La fragmentación internacional de la producción consiste en trasladar las fases de la cadena de valor con mayor contenido tecnológico a países con salarios bajos.} \\
    \textbf{Falsa.} \textit{Justificación: La página 24 indica que la deslocalización consiste en el traslado de ``las fases de la cadena de valor más rutinarias e intensivas en trabajo (fases de fabricación del producto)'', no las de mayor contenido tecnológico.}

    \item \textbf{El fenómeno del retorno de la actividad industrial a los países de origen se ve impulsado por el abaratamiento del transporte internacional.} \\
    \textbf{Falsa.} \textit{Justificación: En la página 25 se menciona que una de las causas que impulsan el retorno o relocalización de la industria es precisamente el ``encarecimiento del transporte'', no su abaratamiento.}

\end{enumerate}

\subsection*{3) ESPECIALIZACIÓN PRODUCTIVA Y COMERCIAL}

\begin{enumerate}[label=\textbf{\arabic*.}]
    \setcounter{enumi}{23}

    \item \textbf{El dinamismo exportador reciente de la industria española se fundamenta principalmente en las ventajas competitivas de sus ramas más tradicionales.} \\
    \textbf{Verdadera.} \textit{Justificación: La página 30 destaca que, en las exportaciones, ``El mayor dinamismo se ha registrado en las ramas más tradicionales [...] dadas las ventajas competitivas de España en éstas''.}

    \item \textbf{Las empresas españolas situadas en sectores de alta intensidad tecnológica suelen estar financiadas mayoritariamente por capital residente nacional.} \\
    \textbf{Falsa.} \textit{Justificación: Según la página 33, en las actividades con mayor requerimiento tecnológico predomina la penetración de capital extranjero (generando más del 50\% del VAB), mientras que el capital residente predomina en las tradicionales.}

    \item \textbf{La especialización de España en actividades de tecnología baja se justifica por la gran dimensión media de sus empresas y su elevada dotación de capital físico.} \\
    \textbf{Falsa.} \textit{Justificación: La página 33 expone lo contrario: las actividades tradicionales se caracterizan por una ``Pequeña dimensión de las empresas, mayor intensidad en recursos naturales y mano de obra'' y ``menor cantidad de capital físico, humano y tecnológico''.}

    \item \textbf{La competencia de países del este de Europa ha mermado la cuota de exportación española en manufacturas de alto contenido tecnológico como las TIC.} \\
    \textbf{Verdadera.} \textit{Justificación: La página 34 lo confirma indicando que la instalación de multinacionales en Hungría, República Checa y Polonia ha superado la cuota española en las exportaciones mundiales TIC.}

\end{enumerate}

\subsection*{4) EFICIENCIA PRODUCTIVA}

\begin{enumerate}[label=\textbf{\arabic*.}]
    \setcounter{enumi}{27}

    \item \textbf{Durante las últimas décadas, la productividad del trabajo en la industria española ha crecido a un ritmo superior al promedio de la Eurozona.} \\
    \textbf{Falsa.} \textit{Justificación: La página 36 señala que, entre 1995 y 2024, la productividad aumentó en España a ``una tasa inferior al promedio de la Eurozona (1,4\% frente al 1,9\%)''.}

    \item \textbf{La mejora de la competitividad en costes lograda durante la Gran Recesión se basó en una devaluación interna mediante la moderación salarial y el ajuste de empleo.} \\
    \textbf{Verdadera.} \textit{Justificación: La página 37 menciona que en la Gran Recesión se mejoró la competitividad en costes debido a la destrucción de empleo, cierre de empresas menos eficientes y moderación salarial, provocando una ``Devaluación interna''.}

    \item \textbf{La fortaleza de las exportaciones industriales españolas se explica porque todo su tejido empresarial presenta niveles de eficiencia productiva muy homogéneos.} \\
    \textbf{Falsa.} \textit{Justificación: La página 41 atribuye la paradoja de las exportaciones a la gran ``Heterogeneidad del tejido empresarial'', donde conviven grandes empresas muy eficientes con otras de menor tamaño y baja eficiencia.}

    \item \textbf{España mantiene su competitividad en ciertos segmentos industriales aprovechando una mano de obra más barata en relación con los países europeos más avanzados.} \\
    \textbf{Verdadera.} \textit{Justificación: La página 41 indica que parte de la especialización productiva se concentra en segmentos de baja productividad, ``aprovechando una mano de obra barata respecto a los países europeos más avanzados''.}

    \item \textbf{La paradoja de las exportaciones españolas reside en que estas crecen a pesar de que la calidad de los productos es inferior a la de sus competidores.} \\
    \textbf{Falsa.} \textit{Justificación: La página 41 resuelve esta paradoja indicando como punto fuerte la ``Mayor calidad comparada de determinadas producciones industriales. Favorable relación calidad-precio'', no una calidad inferior.}

\end{enumerate}

\subsection*{5) POLÍTICA INDUSTRIAL}

\begin{enumerate}[label=\textbf{\arabic*.}]
    \setcounter{enumi}{32}

    \item \textbf{Históricamente, en España ha predominado una visión liberal de la política industrial que desconfiaba de la intervención pública directa.} \\
    \textbf{Verdadera.} \textit{Justificación: La página 43 afirma literalmente que ``En España ha prevalecido una política industrial de corte liberal, que desconfiaba de la viabilidad de la intervención pública''.}

    \item \textbf{La creación de la SEPI en 2001 tuvo como finalidad principal revertir las privatizaciones de los años 90 para recuperar el control estatal de la industria.} \\
    \textbf{Falsa.} \textit{Justificación: La página 43 explica que la finalidad de la SEPI en 2001 fue ``Reordenar el sector público empresarial'' integrando las empresas públicas restantes, no para revertir el proceso de privatización.}

    \item \textbf{La política industrial española actual se define como ``vertical'', ya que se centra exclusivamente en proteger ramas específicas mediante subvenciones directas.} \\
    \textbf{Falsa.} \textit{Justificación: La página 43 especifica que la nueva política industrial se concibe de ``naturaleza horizontal'', dirigiéndose a mejorar las condiciones generales (innovación, internacionalización) para impulsar la productividad, no en proteger ramas verticales específicas.}

\end{enumerate}
